% ============================================================
%  CHAPTER: RELATED WORK
%  Moved from the background chapter into its own chapter.
%  Contains related work for both Operation Mango and Artiphishell.
% ============================================================
\chapter{Related Work}
\label{ch:related-work}


% ============================================================
%  Related Work — Operation Mango
% ============================================================

\section{Operation Mango}
\label{sec:related}

\noindent
\textbf{Static solutions.}
\karonte~\cite{karonte} was one of the first to introduce multi-binary interactions to static analysis in the context of IoT devices.
Their approach models several input sources, such as environment variables and NVRAM entries, and propagates that information from producer binaries to consumer binaries.
\satc~\cite{SaTC} continues their approach, but instead of including network interactions as their sources, they use keywords found in front-end files such as JavaScript and PHP.
They capture directly exploitable vulnerabilities via user input in these front-end files.
This increases the likelihood that \trupoc{}s will be true positives.
However, they ignore other critical binaries in favor of their border binaries.
Both \karonte and \satc severely limit the scope of analysis for efficiency and reduced false positive rates.
For example, \karonte and \satc only analyze binaries that are obviously reachable from the outside network; \satc limits input sources even more by enforcing keyword matches.
As such, they both suffer from high numbers of missed vulnerabilities, as is evidenced by the discrepancy of true positives between \mango and \satc on \karonte's dataset.
\mango avoids this limitation by significantly increasing the analysis efficiency through design benefits (using static data-flow analysis instead of taint tracking and symbolic exploration) and algorithmic improvements (\rev and \assumed).

Saluki~\cite{gotovchits2018saluki} produces an extremely fast static analysis through the micro execution of paths in the binary.
Though this approach is much faster than previous work, it requires precise vulnerability descriptions using their custom language, making the tool difficult to use and highly specific.
\mango takes a much more general approach where categories of sinks defined by function names are set.
Arbiter~\cite{vadayath2022arbiter} begins to cross into the boundary of dynamic analysis with their use of Under-Constrained Symbolic Execution (UCSE).
They approach the problem of vulnerability verification with a static taint analysis followed by UCSE for reachability verification.
However, their approach only works on x86-64 binaries, unlike \mango, which also runs on MIPS and ARM binaries.
Firmalice~\cite{shoshitaishvili2015firmalice} provides backdoor detection and verification in firmware binaries using symbolic execution but is bound by the overhead of symbolic exploration.
Similarly, FIE~\cite{davidson2013fie} uses symbolic execution to analyze open-source MSP430 firmware binaries, but there are various firmware binaries that are intractable to analyze with this tool.

\medskip
\noindent
\textbf{Firmware rehosting and dynamic analysis.}
Firmware rehosting refers to the process of running or emulating firmware on workstations or servers for interaction and security auditing.
Most work in this area conducts full-system emulation of firmware targets, but several have extracted pieces to be analyzed separately~\cite{redini2017bootstomp,costin2014large}.
Once rehosting succeeds, researchers usually move on to vulnerability testing (e.g., using Metasploit) or fuzzing~\cite{chen2018iotfuzzer,ruge2020frankenstein}.
Costin et al.~\cite{costin2014large} used full-system emulation to rehost COTS firmware, and then used web penetration tools to find vulnerabilities in web servers on rehosted firmware.
Firmadyne~\cite{chen2016towards}, Avatar~\cite{zaddach2014avatar}, Avatar$^2$~\cite{muench2018avatar2}, and FirmAE~\cite{kim2020firmae} provides more automation into the full-system rehosting process by QEMU-system to rehost firmware images in a large scale.

Unfortunately, rehosting is usually exceptionally tedious and error-prone, and access to peripherals only available on the original IoT devices is lost.
Therefore, PROSPECT~\cite{kammerstetter2014prospect}, CHARM~\cite{talebi2018charm}, and SURROGATES~\cite{koscher2015surrogates} conducts hardware-in-the-loop operations during firmware emulation.
HALucinator~\cite{clements2020halucinator} and DICE~\cite{mera2021dice} emulate Hardware Abstraction Layers (HALs) and the DMA (Direct Memory Access) channel.
Usually, fuzzing is a natural application after firmware rehosting~\cite{zheng2019firm,zheng2022efficient,srivastava2019firmfuzz,feng2021snipuzz}.


% ============================================================
%  Related Work — Artiphishell
% ============================================================

\section{Artiphishell}
\label{sec:rw-artiphishell}

\noindent
\textbf{Cyber reasoning systems.}
The Cyber Grand Challenge (CGC)~\cite{darpa2016cgc} produced the first generation of autonomous CRSs such as Mechanical Phish~\cite{shoshitaishvili2018mechanical}, which coordinated multiple analysis backends under real-time competition constraints.
These systems operated on simplified binaries under the DECREE OS with fixed system call interfaces.
AIxCC~\cite{aixcc} extends this paradigm to real-world open-source C, C++, and Java software with full build systems, complex dependencies, and realistic vulnerability classes.
Zhang et al.~\cite{zhang2026sokaixcc} provide a systematization of AIxCC finalist architectures and the competition's lessons for autonomous vulnerability discovery.

\medskip
\noindent
\textbf{Coverage-guided and grammar-based fuzzing.}
AFL++~\cite{aflplusplus} is the predominant coverage-guided greybox fuzzer, integrating power schedules~\cite{bohme2016aflfast}, mutation optimization~\cite{lyu2019mopt}, and input-to-state feedback~\cite{aschermann2019redqueen} into a modular framework.
For structured inputs, grammar-based fuzzers such as Nautilus~\cite{nautilus} and Grammarinator~\cite{grammarinator} generate syntactically valid test cases that penetrate past parsing stages.
Jazzer~\cite{jazzer} is a java fuzzer with bytecode-level instrumentation in the JVM.
IJON~\cite{aschermann2020ijon} explores deep state spaces via developer-annotated program states.
LLMs have also been applied to fuzzing: TitanFuzz~\cite{deng2023titanfuzz} generates API call sequences without a seed corpus, ChatAFL~\cite{meng2024chatafl} uses LLM-guided protocol message generation to reach deeper server states, and Fuzz4All~\cite{xia2024fuzz4all} generalizes LLM-based input generation across programming languages.

\medskip
\noindent
\textbf{Static vulnerability detection.}
Pattern-based static analysis tools such as CodeQL~\cite{codeql} and Semgrep~\cite{semgrep} encode vulnerability patterns as declarative queries, trading depth for speed.
Learning-based approaches predict vulnerable components from call-graph features~\cite{Neuhaus2007}, complexity metrics~\cite{Moshtari2013}, or transformer-based source and binary representations~\cite{thapa2022vuldeepecker}.
Bran~\cite{meng2021:Bran} narrows the search space in large repositories by learning syntactic and semantic bug symptoms.

\medskip
\noindent
\textbf{Automated program repair.}
GenProg~\cite{legoues2012genprog} pioneered generate-and-validate repair via genetic programming; SemFix~\cite{nguyen2013semfix} introduced semantic repair through symbolic path conditions; and Angelix~\cite{mechtaev2016angelix} scaled semantic repair to multiline patches.
These classical techniques depend on high-quality test suites for validation.
LLM-based repair does not require as much available information to make repairs, Pearce et al.~\cite{pearce2023zeroshotvulnrepair} showed that LLMs can produce valid patches for certain CWE classes without fine-tuning, and PatchAgent~\cite{yu2025patchagent} decomposes repair into a multi-stage workflow (analysis, root cause localization, targeted synthesis) that mirrors expert reasoning.

\medskip
\noindent
\textbf{Root cause analysis.}
AURORA~\cite{aurora} localizes root causes by statistically comparing execution traces of crashing and non-crashing inputs.
BENZENE~\cite{benzene} mutates program state under under-constrained conditions to isolate causally responsible memory operations.
